% Constructed and checked from the problem; not a geometric proof.
\documentclass[tikz,border=8pt]{standalone}
\begin{document}
\begin{tikzpicture}[scale=1,line cap=round,line join=round]
\coordinate (p0) at (0,5);
\coordinate (p1) at (0,0);
\coordinate (p2) at (7,0);
\coordinate (p3) at (3.5,2.5);
\coordinate (p4) at (3.5,0);
\coordinate (p5) at (0,2.5);
\coordinate (p6) at (7,5);
\coordinate (p7) at (3.5,-2.5);
\draw[thick] (p1) -- (p0);
\draw[thick] (p2) -- (p1);
\draw[thick] (p2) -- (p0);
\draw[dashed] (p5) -- (p3);
\draw[dashed] (p1) -- (p6);
\draw[dashed] (p3) -- (p7);
\draw[dashed] (p1) -- (p5);
\draw[dashed] (p3) -- (p4);
\draw[dashed] (p4) -- (p1);
\draw[dashed] (p0) -- (p6);
\draw[dashed] (p6) -- (p2);
\draw[dashed] (p1) -- (p7);
\draw[dashed] (p7) -- (p2);
\draw[dashed] (p2) -- (p3);
\draw[dashed] (p3) -- (p1);
\draw (1.8191673,3.8468343) -- (1.6808327,3.6531657);
\draw (5.3191673,1.3468343) -- (5.1808327,1.1531657);
\draw (1.70835,0.119) -- (1.70835,-0.119);
\draw (1.79165,0.119) -- (1.79165,-0.119);
\draw (5.20835,0.119) -- (5.20835,-0.119);
\draw (5.29165,0.119) -- (5.29165,-0.119);
\fill (p0) circle (1.1pt) node[fill=white,inner sep=1pt,above left] {$A$};
\fill (p1) circle (1.1pt) node[fill=white,inner sep=1pt,below left] {$B$};
\fill (p2) circle (1.1pt) node[fill=white,inner sep=1pt,below right] {$C$};
\fill (p3) circle (1.1pt) node[fill=white,inner sep=1pt,above right] {$E$};
\fill (p4) circle (1.1pt) node[fill=white,inner sep=1pt,below right] {$F$};
\fill (p5) circle (1.1pt) node[fill=white,inner sep=1pt,above left] {$M$};
\fill (p6) circle (1.1pt) node[fill=white,inner sep=1pt,above right] {$K$};
\fill (p7) circle (1.1pt) node[fill=white,inner sep=1pt,below right] {$G$};
\end{tikzpicture}
\end{document}
